%! Mode:: "TeX:UTF-8"
%! TEX program = xelatex
\documentclass[withoutpreface,bwprint]{cumcmthesis}
\usepackage{etoolbox}
\BeforeBeginEnvironment{tabular}{\zihao{-5}}
\usepackage[numbers,sort&compress]{natbib}
\usepackage[framemethod=TikZ]{mdframed}
\usepackage{float}
\usepackage[ruled,linesnumbered]{algorithm2e}
\AtBeginEnvironment{algorithm}{\normalem}
\renewcommand{\listalgorithmcfname}{算法列表}
\renewcommand{\algorithmcfname}{算法}
\renewcommand{\algorithmautorefname}{算法}
\renewcommand{\algorithmcflinename}{行}
\crefname{algorithm}{算法}{算法}
\crefname{algocf}{算法}{算法}
\Crefname{algocf}{算法}{算法}
\usepackage{tikz}
\usetikzlibrary{arrows.meta,positioning,shapes.geometric,calc,fit,backgrounds}
\usepackage{colortbl}
\usepackage{booktabs}
\usepackage{multirow}
\usepackage{array}
\usepackage{longtable}
\usepackage{pifont}
\usepackage{pdflscape}

\newcommand{\xiaosan}{\fontsize{15pt}{18pt}\selectfont}

% mcmicm_styles.tex
% 从美赛模板提取的代码/AI 报告样式
% 使用时导言区加入: % mcmicm_styles.tex
% 从美赛模板提取的代码/AI 报告样式
% 使用时导言区加入: % mcmicm_styles.tex
% 从美赛模板提取的代码/AI 报告样式
% 使用时导言区加入: \input{mcmicm_styles.tex}

\usepackage{tcolorbox} % 美丽盒子包
\usepackage{matlab-prettifier} % MATLAB代码排版包
\tcbuselibrary{listings,skins,breakable,xparse}

% 颜色定义
\definecolor{darkgray}{rgb}{0.4,0.4,0.4}
\definecolor{darkblue}{rgb}{0.0,0.0,0.5}
\definecolor{darkgreen}{rgb}{0,0.45,0.13}
\definecolor{orange}{rgb}{0.87,0.49,0.0}

% ========== tcolorbox 列表样式 ==========
\tcbset{
    mylist/.style = {
        colframe = gray,
        colback = white,
        coltitle = red!50!yellow!3!white,
        colbacktitle = white,
        listing only,
        attach boxed title to top center = {yshift = -\tcboxedtitleheight/2},
        enhanced,
        drop fuzzy shadow,
        left = 6.5mm,
        breakable,
        fonttitle=\small\bfseries\color{black},
        before skip=12pt plus 2pt,
        after skip=12pt plus 2pt,
    },
    example/.style 2 args = {
        mylist,
        title = {Listing \thetcbcounter: #1},
        label = {#2},
    },
}

\newtcblisting[auto counter, number within = section, list inside = mcode]{matlab}[3][]{%
    listing options = {
        style = Matlab-editor,
        numbers = left,
        numberstyle = \footnotesize\color{darkgray}\ttfamily,
        basicstyle = \footnotesize\ttfamily,
        breaklines = true,
        breakatwhitespace = true,
        postbreak = \mbox{\textcolor{gray}{$\hookrightarrow$}\space}
    },
    overlay = {\begin{tcbclipinterior}\fill[blue!15!white] (frame.south west) rectangle ([xshift=5.3mm]frame.north west);\end{tcbclipinterior}},
    example = {#2}{#3}, #1,
}

\newtcblisting[auto counter, number within = section, list inside = pcode]{python}[3][]{%
    listing options = {
        language = Python,
        numbers = left,
        numberstyle = \footnotesize\color{darkgray}\ttfamily,
        basicstyle = \footnotesize\ttfamily,
        keywordstyle = \color{darkblue}\bfseries,
        commentstyle = \color{darkgreen},
        stringstyle = \color{red!60},
        showstringspaces = false,
        emphstyle = \color{orange!60},
        emph = {self, cls, True, False, None},
        breaklines = true,
        breakatwhitespace = true,
        postbreak = \mbox{\textcolor{gray}{$\hookrightarrow$}\space}
    },
    overlay = {\begin{tcbclipinterior}\fill[blue!15!white] (frame.south west) rectangle ([xshift=5.3mm]frame.north west);\end{tcbclipinterior}},
    example = {#2}{#3}, #1,
}

% ========== AI 报告专用无边框 listings 风格 ==========
\lstdefinestyle{ai_python_style}{
    language = Python,
    numbers = none,
    basicstyle = \footnotesize\ttfamily,
    keywordstyle = \color{darkblue}\bfseries,
    commentstyle = \color{darkgreen},
    stringstyle = \color{red!60},
    showstringspaces = false,
    emphstyle = \color{orange!60},
    emph = {self, cls, True, False, None},
    breaklines = true,
    breakatwhitespace = true,
    postbreak = \mbox{\textcolor{gray}{$\hookrightarrow$}\space},
    frame = none,
    backgroundcolor = \color{white},
    xleftmargin = 0pt,
    xrightmargin = 0pt,
}

\lstdefinestyle{ai_matlab_style}{
    style = Matlab-editor,
    numbers = none,
    basicstyle = \footnotesize\ttfamily,
    breaklines = true,
    breakatwhitespace = true,
    postbreak = \mbox{\textcolor{gray}{$\hookrightarrow$}\space},
    frame = none,
    backgroundcolor = \color{white},
    xleftmargin = 0pt,
    xrightmargin = 0pt,
}

\lstnewenvironment{ai_python}[1][]
    {\lstset{style=ai_python_style, #1}}
    {}

\lstnewenvironment{ai_matlab}[1][]
    {\lstset{style=ai_matlab_style, #1}}
    {}

% ========== 为定理类环境添加底色框 ==========
\mdfdefinestyle{thmshaded}{%
    backgroundcolor=gray!8,% 浅灰底色
    linecolor=black!30,% 边框颜色
    leftline=true,% 左侧线条
    topline=false,
    rightline=false,
    bottomline=false,
    linewidth=3pt,% 左侧线宽
    innerleftmargin=10pt,
    innertopmargin=6pt,
    innerbottommargin=6pt,
    roundcorner=0pt,
}

\surroundwithmdframed[style=thmshaded]{definition}
\surroundwithmdframed[style=thmshaded]{theorem}
\surroundwithmdframed[style=thmshaded]{lemma}
\surroundwithmdframed[style=thmshaded]{corollary}
\surroundwithmdframed[style=thmshaded]{assumption}
\surroundwithmdframed[style=thmshaded]{example}
\surroundwithmdframed[style=thmshaded]{problem}
\surroundwithmdframed[style=thmshaded]{solution}




\usepackage{tcolorbox} % 美丽盒子包
\usepackage{matlab-prettifier} % MATLAB代码排版包
\tcbuselibrary{listings,skins,breakable,xparse}

% 颜色定义
\definecolor{darkgray}{rgb}{0.4,0.4,0.4}
\definecolor{darkblue}{rgb}{0.0,0.0,0.5}
\definecolor{darkgreen}{rgb}{0,0.45,0.13}
\definecolor{orange}{rgb}{0.87,0.49,0.0}

% ========== tcolorbox 列表样式 ==========
\tcbset{
    mylist/.style = {
        colframe = gray,
        colback = white,
        coltitle = red!50!yellow!3!white,
        colbacktitle = white,
        listing only,
        attach boxed title to top center = {yshift = -\tcboxedtitleheight/2},
        enhanced,
        drop fuzzy shadow,
        left = 6.5mm,
        breakable,
        fonttitle=\small\bfseries\color{black},
        before skip=12pt plus 2pt,
        after skip=12pt plus 2pt,
    },
    example/.style 2 args = {
        mylist,
        title = {Listing \thetcbcounter: #1},
        label = {#2},
    },
}

\newtcblisting[auto counter, number within = section, list inside = mcode]{matlab}[3][]{%
    listing options = {
        style = Matlab-editor,
        numbers = left,
        numberstyle = \footnotesize\color{darkgray}\ttfamily,
        basicstyle = \footnotesize\ttfamily,
        breaklines = true,
        breakatwhitespace = true,
        postbreak = \mbox{\textcolor{gray}{$\hookrightarrow$}\space}
    },
    overlay = {\begin{tcbclipinterior}\fill[blue!15!white] (frame.south west) rectangle ([xshift=5.3mm]frame.north west);\end{tcbclipinterior}},
    example = {#2}{#3}, #1,
}

\newtcblisting[auto counter, number within = section, list inside = pcode]{python}[3][]{%
    listing options = {
        language = Python,
        numbers = left,
        numberstyle = \footnotesize\color{darkgray}\ttfamily,
        basicstyle = \footnotesize\ttfamily,
        keywordstyle = \color{darkblue}\bfseries,
        commentstyle = \color{darkgreen},
        stringstyle = \color{red!60},
        showstringspaces = false,
        emphstyle = \color{orange!60},
        emph = {self, cls, True, False, None},
        breaklines = true,
        breakatwhitespace = true,
        postbreak = \mbox{\textcolor{gray}{$\hookrightarrow$}\space}
    },
    overlay = {\begin{tcbclipinterior}\fill[blue!15!white] (frame.south west) rectangle ([xshift=5.3mm]frame.north west);\end{tcbclipinterior}},
    example = {#2}{#3}, #1,
}

% ========== AI 报告专用无边框 listings 风格 ==========
\lstdefinestyle{ai_python_style}{
    language = Python,
    numbers = none,
    basicstyle = \footnotesize\ttfamily,
    keywordstyle = \color{darkblue}\bfseries,
    commentstyle = \color{darkgreen},
    stringstyle = \color{red!60},
    showstringspaces = false,
    emphstyle = \color{orange!60},
    emph = {self, cls, True, False, None},
    breaklines = true,
    breakatwhitespace = true,
    postbreak = \mbox{\textcolor{gray}{$\hookrightarrow$}\space},
    frame = none,
    backgroundcolor = \color{white},
    xleftmargin = 0pt,
    xrightmargin = 0pt,
}

\lstdefinestyle{ai_matlab_style}{
    style = Matlab-editor,
    numbers = none,
    basicstyle = \footnotesize\ttfamily,
    breaklines = true,
    breakatwhitespace = true,
    postbreak = \mbox{\textcolor{gray}{$\hookrightarrow$}\space},
    frame = none,
    backgroundcolor = \color{white},
    xleftmargin = 0pt,
    xrightmargin = 0pt,
}

\lstnewenvironment{ai_python}[1][]
    {\lstset{style=ai_python_style, #1}}
    {}

\lstnewenvironment{ai_matlab}[1][]
    {\lstset{style=ai_matlab_style, #1}}
    {}

% ========== 为定理类环境添加底色框 ==========
\mdfdefinestyle{thmshaded}{%
    backgroundcolor=gray!8,% 浅灰底色
    linecolor=black!30,% 边框颜色
    leftline=true,% 左侧线条
    topline=false,
    rightline=false,
    bottomline=false,
    linewidth=3pt,% 左侧线宽
    innerleftmargin=10pt,
    innertopmargin=6pt,
    innerbottommargin=6pt,
    roundcorner=0pt,
}

\surroundwithmdframed[style=thmshaded]{definition}
\surroundwithmdframed[style=thmshaded]{theorem}
\surroundwithmdframed[style=thmshaded]{lemma}
\surroundwithmdframed[style=thmshaded]{corollary}
\surroundwithmdframed[style=thmshaded]{assumption}
\surroundwithmdframed[style=thmshaded]{example}
\surroundwithmdframed[style=thmshaded]{problem}
\surroundwithmdframed[style=thmshaded]{solution}




\usepackage{tcolorbox} % 美丽盒子包
\usepackage{matlab-prettifier} % MATLAB代码排版包
\tcbuselibrary{listings,skins,breakable,xparse}

% 颜色定义
\definecolor{darkgray}{rgb}{0.4,0.4,0.4}
\definecolor{darkblue}{rgb}{0.0,0.0,0.5}
\definecolor{darkgreen}{rgb}{0,0.45,0.13}
\definecolor{orange}{rgb}{0.87,0.49,0.0}

% ========== tcolorbox 列表样式 ==========
\tcbset{
    mylist/.style = {
        colframe = gray,
        colback = white,
        coltitle = red!50!yellow!3!white,
        colbacktitle = white,
        listing only,
        attach boxed title to top center = {yshift = -\tcboxedtitleheight/2},
        enhanced,
        drop fuzzy shadow,
        left = 6.5mm,
        breakable,
        fonttitle=\small\bfseries\color{black},
        before skip=12pt plus 2pt,
        after skip=12pt plus 2pt,
    },
    example/.style 2 args = {
        mylist,
        title = {Listing \thetcbcounter: #1},
        label = {#2},
    },
}

\newtcblisting[auto counter, number within = section, list inside = mcode]{matlab}[3][]{%
    listing options = {
        style = Matlab-editor,
        numbers = left,
        numberstyle = \footnotesize\color{darkgray}\ttfamily,
        basicstyle = \footnotesize\ttfamily,
        breaklines = true,
        breakatwhitespace = true,
        postbreak = \mbox{\textcolor{gray}{$\hookrightarrow$}\space}
    },
    overlay = {\begin{tcbclipinterior}\fill[blue!15!white] (frame.south west) rectangle ([xshift=5.3mm]frame.north west);\end{tcbclipinterior}},
    example = {#2}{#3}, #1,
}

\newtcblisting[auto counter, number within = section, list inside = pcode]{python}[3][]{%
    listing options = {
        language = Python,
        numbers = left,
        numberstyle = \footnotesize\color{darkgray}\ttfamily,
        basicstyle = \footnotesize\ttfamily,
        keywordstyle = \color{darkblue}\bfseries,
        commentstyle = \color{darkgreen},
        stringstyle = \color{red!60},
        showstringspaces = false,
        emphstyle = \color{orange!60},
        emph = {self, cls, True, False, None},
        breaklines = true,
        breakatwhitespace = true,
        postbreak = \mbox{\textcolor{gray}{$\hookrightarrow$}\space}
    },
    overlay = {\begin{tcbclipinterior}\fill[blue!15!white] (frame.south west) rectangle ([xshift=5.3mm]frame.north west);\end{tcbclipinterior}},
    example = {#2}{#3}, #1,
}

% ========== AI 报告专用无边框 listings 风格 ==========
\lstdefinestyle{ai_python_style}{
    language = Python,
    numbers = none,
    basicstyle = \footnotesize\ttfamily,
    keywordstyle = \color{darkblue}\bfseries,
    commentstyle = \color{darkgreen},
    stringstyle = \color{red!60},
    showstringspaces = false,
    emphstyle = \color{orange!60},
    emph = {self, cls, True, False, None},
    breaklines = true,
    breakatwhitespace = true,
    postbreak = \mbox{\textcolor{gray}{$\hookrightarrow$}\space},
    frame = none,
    backgroundcolor = \color{white},
    xleftmargin = 0pt,
    xrightmargin = 0pt,
}

\lstdefinestyle{ai_matlab_style}{
    style = Matlab-editor,
    numbers = none,
    basicstyle = \footnotesize\ttfamily,
    breaklines = true,
    breakatwhitespace = true,
    postbreak = \mbox{\textcolor{gray}{$\hookrightarrow$}\space},
    frame = none,
    backgroundcolor = \color{white},
    xleftmargin = 0pt,
    xrightmargin = 0pt,
}

\lstnewenvironment{ai_python}[1][]
    {\lstset{style=ai_python_style, #1}}
    {}

\lstnewenvironment{ai_matlab}[1][]
    {\lstset{style=ai_matlab_style, #1}}
    {}

% ========== 为定理类环境添加底色框 ==========
\mdfdefinestyle{thmshaded}{%
    backgroundcolor=gray!8,% 浅灰底色
    linecolor=black!30,% 边框颜色
    leftline=true,% 左侧线条
    topline=false,
    rightline=false,
    bottomline=false,
    linewidth=3pt,% 左侧线宽
    innerleftmargin=10pt,
    innertopmargin=6pt,
    innerbottommargin=6pt,
    roundcorner=0pt,
}

\surroundwithmdframed[style=thmshaded]{definition}
\surroundwithmdframed[style=thmshaded]{theorem}
\surroundwithmdframed[style=thmshaded]{lemma}
\surroundwithmdframed[style=thmshaded]{corollary}
\surroundwithmdframed[style=thmshaded]{assumption}
\surroundwithmdframed[style=thmshaded]{example}
\surroundwithmdframed[style=thmshaded]{problem}
\surroundwithmdframed[style=thmshaded]{solution}




%=============================================================================
% 自定义样式
%=============================================================================
\definecolor{clkSys}{RGB}{214,234,248}
\definecolor{clkPix}{RGB}{255,235,205}
\definecolor{clkHdmi}{RGB}{212,240,218}
\definecolor{boxFill}{RGB}{245,247,250}
\definecolor{boxFrame}{RGB}{70,100,140}

\tikzset{
    blk/.style   = {draw=boxFrame, fill=boxFill, rounded corners=2pt,
                    minimum width=2.6cm, minimum height=0.72cm,
                    align=center, font=\zihao{-5}},
    blkS/.style  = {blk, minimum width=1.9cm},
    blkWide/.style = {blk, minimum width=4.2cm},
    io/.style    = {draw=boxFrame!70, fill=white, rounded corners=6pt,
                    minimum width=2.2cm, minimum height=0.9cm,
                    align=center, font=\zihao{-5}},
    ar/.style    = {-{Latex[length=1.6mm]}, draw=boxFrame, semithick},
    arD/.style   = {{Latex[length=1.6mm]}-{Latex[length=1.6mm]},
                    draw=boxFrame, semithick},
}

\title{音速并行：基于 Artix-7 的纯 RTL 音频一体化系统\\设计报告}
\author{Binyu Fei}
\date{2026 年 10 月}

%%%%%%%%%%%%%%%%%%%%%%%%%%%%%%%%%%%%%%%%%%%%%%%%%%%%%%%%%%%%%%%%%%%%%%%%%%%%%%%
\begin{document}
%%%%%%%%%%%%%%%%%%%%%%%%%%%%%%%%%%%%%%%%%%%%%%%%%%%%%%%%%%%%%%%%%%%%%%%%%%%%%%%
\maketitle

\begin{abstract}
本作品在 Xilinx Artix-7 XC7A100T 上以\textbf{纯 RTL}（不使用 HLS、软核或任何现成 IP）
实现了一套音频一体化系统，覆盖“采集—均衡—频谱分析—可视化—交互”完整闭环。
系统以 WM8960 为音频前端（48\,kHz/24\,bit），采样数据同时进入两条并行流水线：
5 段定点 IIR 参量均衡器，以及 1024 点基 2 FFT 频谱分析引擎。FFT 仅用\textbf{一个}
蝶形单元配合双口 BRAM 时分复用，硬件开销不到全并行方案的 1/512，整帧处理仅占
帧周期的 1.3\%。频谱结果经对数分箱、峰值保持与线性衰减后驱动\textbf{双路独立显示}：
LCD 480$\times$272 作交互界面（8 按钮面板 + 8 行状态栏），HDMI 640$\times$360 作展示大屏
（80 根独立柱 + 波形），并从 SD 卡读取图片作为动态背景。
系统横跨四个时钟域，以格雷码异步 FIFO 与电平翻转握手完成跨域传递。
整机占用 11\,716 LUT（18.48\%）、100.5 BRAM、26 DSP48，时序全部收敛，
\textbf{综合与实现 0 错误 0 警告}，41 个测试平台全部与 Python 定点黄金模型逐位对拍通过。

\keywords{FPGA \quad 纯 RTL \quad 定点 DSP \quad 频谱可视化 \quad HDMI \quad 跨时钟域}
\end{abstract}

\tableofcontents
\newpage

%%%%%%%%%%%%%%%%%%%%%%%%%%%%%%%%%%%%%%%%%%%%%%%%%%%%%%%%%%%%%%%%%%%%%%%%%%%%%%%
\section{选题背景与创新点}
%%%%%%%%%%%%%%%%%%%%%%%%%%%%%%%%%%%%%%%%%%%%%%%%%%%%%%%%%%%%%%%%%%%%%%%%%%%%%%%

\subsection{为什么选这个题目}

音频可视化设备在消费市场与工业检测中需求广泛，但现有方案存在结构性矛盾，
如表~\ref{tab:bg-compare} 所示。

\begin{table}[H]
\centering
\caption{现有方案的结构性矛盾}\label{tab:bg-compare}
\begin{tabular}{p{4.2cm}p{9.5cm}}
\toprule[1.5pt]
\rowcolor{gray!25} 方案 & 问题 \\
\midrule[1pt]
PC + 声卡 + 上位机 & 延迟高（10\textasciitilde 100\,ms），功耗大，不适合嵌入式部署 \\
MCU + CMSIS-DSP   & 算力不足，FFT + EQ + 双屏驱动无法同时满足 \\
Zynq + PYNQ       & 依赖 HLS/Python 框架，底层时序不可控，启动秒级 \\
商用 ASIC         & 功能固定，无法在线调参，不便于教学研究 \\
\bottomrule[1.5pt]
\end{tabular}
\end{table}

\textbf{核心矛盾}是“实时性”与“可调性”难以兼得：专用 ASIC 实时但不灵活，
软件方案灵活但延迟大。本作品用纯 RTL 在低端 FPGA 上实现一条
\textbf{端到端延迟可预测}的音频处理链路，做到确定性延迟、在线可调、零 IP 依赖。

\subsection{创新点}

\subsubsection{创新点一：扫描式双路渲染，无帧缓冲}

传统做法先渲染到 framebuffer 再扫描读出，一帧 640$\times$360$\times$16\,bit
需 103 块 BRAM，而本器件只有 135 块。本作品把渲染核心抽象为纯组合函数
$(x,y,\text{mode})\rightarrow$\,RGB，\textbf{扫描到哪个像素就算哪个像素}，
一份参数化的渲染核心实例化两次即可驱动两种分辨率。

\subsubsection{创新点二：单蝶形 FFT 的极低资源占用}

1024 点 FFT 若全并行需 512 个蝶形 $\times$ 4 乘法器 $=$ 2048 个 DSP48，
超出器件总量（240 个）8 倍。本作品用\textbf{一个}蝶形单元加双口 BRAM
in-place 时分复用，DSP 占用降到 4 个，整帧耗时仅占帧周期 1.3\%，余量 75 倍。

\subsubsection{创新点三：数据柱数与显示柱数解耦}

LCD 480 像素宽、60 根柱（每柱 8 像素，整除）；HDMI 640 像素宽，若 60 根柱
则每柱 10.67 像素（不整除，需乘法）。本作品把“有多少数据”与“显示多少根柱”
拆成两个参数（\texttt{NBARS} / \texttt{NBARS\_DISP}），并用 \texttt{generate}
的\textbf{恒等分支}保证老路径在综合后结构完全不变。

\subsubsection{创新点四：不同尺寸的双帧缓冲替代“双读口”}

两个显示器的时钟不同（12.5\,MHz / 25\,MHz）。实测表明 Vivado 的 BRAM 推断
\textbf{不支持“1 写 2 读跨时钟”}，会静默退化成分步式 RAM（LUTRAM），
触发 \texttt{DRC UTLZ-1}（需 86\,452 个 \texttt{RAMD64E}，仅 19\,000 个可用），
加 \texttt{(* ram\_style = "block" *)} 亦无效。本作品改为\textbf{两份不同尺寸的缓冲}：
LCD 直读 480$\times$272 全图，HDMI 读 320$\times$180 裁剪区后 2 倍整数放大，
BRAM 从 116 块降到 84 块。

\subsubsection{创新点五：三级知识沉淀的开发方法}

问题账本（825 条）$\rightarrow$ 技能包（4 个）$\rightarrow$ 状态快照，
把每次失败转化为可复用的方法论，详见第~\ref{sec:skill} 章。

%%%%%%%%%%%%%%%%%%%%%%%%%%%%%%%%%%%%%%%%%%%%%%%%%%%%%%%%%%%%%%%%%%%%%%%%%%%%%%%
\section{设计原理与功能框图}
%%%%%%%%%%%%%%%%%%%%%%%%%%%%%%%%%%%%%%%%%%%%%%%%%%%%%%%%%%%%%%%%%%%%%%%%%%%%%%%

\subsection{系统功能框图}

%=============================================================================
% fig_arch.tex —— 模块级系统框图（由 main.tex \input）
%   风格参考：PL/PS 分区 + 功能分组灰底框 + 每模块一个小方块
%=============================================================================

\begin{landscape}
\begin{figure}[H]
\centering
\resizebox{\linewidth}{!}{%
\begin{tikzpicture}[
    font=\zihao{6},
    mod/.style   = {draw=black!70, fill=white, minimum width=1.55cm,
                    minimum height=0.50cm, align=center, inner sep=1pt,
                    rounded corners=1pt},
    modW/.style  = {mod, minimum width=2.35cm},
    mem/.style   = {draw=black!55, fill=gray!12, minimum width=1.30cm,
                    minimum height=1.10cm, align=center, inner sep=1pt},
    io/.style    = {draw=black!75, fill=gray!22, minimum width=1.95cm,
                    minimum height=1.90cm, align=center, inner sep=2pt,
                    rounded corners=2pt},
    ioSm/.style  = {draw=black!75, fill=gray!22, minimum width=1.95cm,
                    minimum height=0.95cm, align=center, inner sep=2pt,
                    rounded corners=2pt},
    grp/.style   = {draw=black!60, fill=black!7, rounded corners=3pt,
                    inner sep=5pt},
    grpB/.style  = {grp, fill=blue!6},
    grpO/.style  = {grp, fill=orange!10},
    grpG/.style  = {grp, fill=green!8},
    grpP/.style  = {grp, fill=violet!7},
    ar/.style    = {-{Latex[length=1.3mm,width=1.0mm]}, draw=black!65,
                    semithick},
    arBus/.style = {-{Latex[length=1.3mm,width=1.0mm]}, draw=black!65,
                    line width=1.6pt},
    arD/.style   = {{Latex[length=1.3mm,width=1.0mm]}-{Latex[length=1.3mm,width=1.0mm]},
                    draw=black!65, semithick},
    lbl/.style   = {font=\zihao{6}, inner sep=1pt, fill=white,
                    fill opacity=0.85, text opacity=1},
    grplbl/.style= {font=\zihao{-6}\bfseries, inner sep=1pt, text=black!75},
]


%=============================================================================
% 图 A：音频采集 → 均衡 → 频谱分析
%=============================================================================
\node[io]   (codec) at (0, -3.0)  {WM8960\\编解码器};
\node[ioSm] (uart)  at (0, -7.4)  {PC 终端\\(CH340E)};

% ---- 时钟与复位 ----
\node[modW] (mmcm1) at (3.4, -0.9) {\texttt{clk\_gen} / MMCM1\\48 + 12.5\,MHz};
\node[modW] (rsts)  at (3.4, -2.0) {\texttt{rst\_sync} $\times$3};

% ---- Codec 配置 ----
\node[modW] (i2cc)  at (3.4, -3.6) {\texttt{i2c\_control}};
\node[modW] (wminit)at (3.4, -4.6) {\texttt{WM8960\_init}};
\node[modW] (wmrom) at (3.4, -5.6) {\texttt{WM8960\_init\_table}};

% ---- 音频链路 ----
\node[modW] (i2sclk) at (7.0, -2.4) {\texttt{i2s\_slave\_clk}};
\node[modW] (i2srx)  at (7.0, -3.4) {\texttt{i2s\_rx}};
\node[modW] (dcblk)  at (7.0, -4.4) {\texttt{dc\_block} $\times$2};
\node[modW] (eq)     at (7.0, -5.4) {\texttt{eq\_cascade} $\times$2\\5 段双二阶};
\node[modW] (eqrom)  at (7.0, -6.6) {\texttt{eq\_coeff\_rom}};
\node[modW] (i2stx)  at (7.0, -7.7) {\texttt{i2s\_tx}};
\node[modW] (afifo)  at (7.0, -8.8) {\texttt{audio\_fifo} $\times$2};

% ---- 频谱分析 ----
\node[modW] (fftctrl)at (10.9, -2.4) {\texttt{fft\_core}\\控制状态机};
\node[modW] (bfly)   at (10.9, -3.4) {\texttt{fft\_butterfly}\\单蝶形时分复用};
\node[modW] (addrgen)at (10.9, -4.4) {\texttt{fft\_addr\_gen}};
\node[modW] (twrom)  at (10.9, -5.4) {\texttt{fft\_twiddle\_rom}};
\node[mem]   (wram)  at (10.9, -7.0) {工作\\RAM\\1024$\times$24};
\node[modW] (mag)    at (10.9, -8.4) {幅度估计\\$\max+\frac{7}{16}\min$};
\node[modW] (spect)  at (10.9, -9.4) {\texttt{spectrum}\\对数分箱+峰值保持};
\node[modW] (smap)   at (10.9, -10.4){\texttt{spectrum\_map}\\\texttt{spectrum\_map80}};

% ---- 跨时钟域 ----
\node[modW] (async)  at (14.6, -2.4) {\texttt{async\_fifo}};
\node[modW] (specsync)at(14.6, -3.4) {\texttt{spec\_sync}};
\node[modW] (wave)   at (14.6, -4.6) {\texttt{wave\_buf}\\触发+双时钟};
\node[modW] (fir)    at (14.6, -5.8) {\texttt{fir\_sym} (15 抽头)};
\node[modW] (env)    at (14.6, -7.0) {\texttt{env\_track}};

% ---- 控制 ----
\node[modW] (kdb)    at (14.6, -8.6) {\texttt{key\_debounce}};
\node[modW] (kgst)   at (14.6, -9.6) {\texttt{key\_gesture}};
\node[modW] (uartrx) at (14.6, -10.6){\texttt{uart\_rx} / \texttt{uart\_tx}};
\node[modW] (cmdpr)  at (14.6, -11.6){\texttt{cmd\_proc}};
\node[modW] (uictrl) at (14.6, -12.6){\texttt{ui\_ctrl}};

\begin{scope}[on background layer]
\node[grpG, fit=(mmcm1)(rsts), label={[grplbl]above:时钟与复位}] (gClk) {};
\node[grpB, fit=(i2cc)(wminit)(wmrom), label={[grplbl]above:Codec 配置}] (gCodec) {};
\node[grpB, fit=(i2sclk)(i2srx)(dcblk)(eq)(eqrom)(i2stx)(afifo),
      label={[grplbl]above:音频链路 (clk\_sys 48\,MHz)}] (gAudio) {};
\node[grpB, fit=(fftctrl)(bfly)(addrgen)(twrom)(wram)(mag)(spect)(smap),
      label={[grplbl]above:频谱分析 (clk\_sys)}] (gSpec) {};
\node[grpG, fit=(async)(specsync)(wave)(fir)(env),
      label={[grplbl]above:跨时钟域 + 波形}] (gCdc) {};
\node[grpP, fit=(kdb)(kgst)(uartrx)(cmdpr)(uictrl),
      label={[grplbl]above:交互控制}] (gCtrl) {};
\node[draw=black!80, line width=1.0pt, rounded corners=6pt,
      fit=(gClk)(gCodec)(gAudio)(gSpec)(gCdc)(gCtrl), inner sep=8pt,
      label={[font=\zihao{-5}\bfseries]above:FPGA —— XC7A100T-2FGG484I}] (fpga) {};
\end{scope}

\draw[arD] (codec.east) ++(0,0.5) -- node[lbl,above]{I2S 从} ([yshift=0.6cm]gAudio.west);
\draw[arD] (codec.east) ++(0,0.1) -- node[lbl,above]{I2C} ([yshift=-0.6cm]gCodec.west);
\draw[ar]  (codec.north) -- node[lbl,right]{MCLK} ([yshift=0.2cm]gClk.south -| codec.north);
\draw[arD] (uart.east) -- node[lbl,above]{115200} ([yshift=0.5cm]gCtrl.west);
\draw[arBus] (gAudio.east) -- node[lbl,above]{rx\_l / rx\_r (Q1.23)} ([yshift=-1.6cm]gSpec.west);
\draw[arBus] (gAudio.east) -- ++(0.35,0) |- node[lbl,pos=0.25,above]{波形抽头} ([yshift=-3.2cm]gCdc.west);
\draw[ar, dashed, draw=blue!55] ([yshift=0.7cm]gClk.east) -- ++(0.4,0) |-
     node[lbl,pos=0.3,above]{\texttt{clk\_sys} 48\,MHz} ([yshift=0.9cm]gAudio.west);
\draw[arBus] (gSpec.east) -- ++(0.5,0) |- node[lbl,pos=0.2,above]{柱高/波形快照} ([yshift=-1.6cm]gCdc.west);

\end{tikzpicture}%
}
\caption{模块级框图（一）：音频采集 $\rightarrow$ 均衡 $\rightarrow$ FFT 频谱分析，及跨时钟域与控制}
\label{fig:arch-a}
\end{figure}
\end{landscape}

\begin{landscape}
\begin{figure}[H]
\centering
\resizebox{\linewidth}{!}{%
\begin{tikzpicture}[
    font=\zihao{6},
    mod/.style   = {draw=black!70, fill=white, minimum width=1.55cm,
                    minimum height=0.50cm, align=center, inner sep=1pt,
                    rounded corners=1pt},
    modW/.style  = {mod, minimum width=2.35cm},
    mem/.style   = {draw=black!55, fill=gray!12, minimum width=1.30cm,
                    minimum height=1.10cm, align=center, inner sep=1pt},
    io/.style    = {draw=black!75, fill=gray!22, minimum width=1.95cm,
                    minimum height=1.90cm, align=center, inner sep=2pt,
                    rounded corners=2pt},
    ioSm/.style  = {draw=black!75, fill=gray!22, minimum width=1.95cm,
                    minimum height=0.95cm, align=center, inner sep=2pt,
                    rounded corners=2pt},
    grp/.style   = {draw=black!60, fill=black!7, rounded corners=3pt, inner sep=5pt},
    grpB/.style  = {grp, fill=blue!6},
    grpO/.style  = {grp, fill=orange!10},
    grpG/.style  = {grp, fill=green!8},
    grpP/.style  = {grp, fill=violet!7},
    ar/.style    = {-{Latex[length=1.3mm,width=1.0mm]}, draw=black!65, semithick},
    arBus/.style = {-{Latex[length=1.3mm,width=1.0mm]}, draw=black!65, line width=1.6pt},
    arD/.style   = {{Latex[length=1.3mm,width=1.0mm]}-{Latex[length=1.3mm,width=1.0mm]},
                    draw=black!65, semithick},
    lbl/.style   = {font=\zihao{6}, inner sep=1pt, fill=white, fill opacity=0.85, text opacity=1},
    grplbl/.style= {font=\zihao{-6}\bfseries, inner sep=1pt, text=black!75},
]

%=============================================================================
% 图 B：显示链 + SD 卡 + HDMI
%=============================================================================
\node[io]   (lcd)   at (0, -2.6)  {LCD\\480$\times$272};
\node[io]   (hdmi)  at (0, -7.0)  {HDMI 显示器\\640$\times$360};
\node[ioSm] (sdcard)at (0, -10.6) {microSD 卡};

% ---- 显示核心（两份实例）----
\node[modW] (timing) at (4.2, -0.9) {\texttt{lcd\_timing}\\12.5M / 25M 参数化};
\node[modW] (bg)     at (4.2, -2.0) {\texttt{bg\_src} (渐变归一化)};
\node[modW] (polar)  at (4.2, -3.1) {\texttt{polar\_map} 64 扇区};
\node[modW] (polrom) at (4.2, -4.2) {\texttt{polar\_map\_rom}};
\node[modW] (mix)    at (4.2, -5.3) {\texttt{disp\_mix}\\背景+柱+波形};
\node[modW] (rbrom)  at (4.2, -6.5) {\texttt{rainbow\_rom}};
\node[modW] (demo)   at (4.2, -7.6) {\texttt{demo\_src}};
\node[modW] (wfilt)  at (4.2, -8.7) {\texttt{wave\_buf} + \texttt{fir\_sym}};

% ---- UI 层 ----
\node[modW] (anim)   at (8.2, -0.9) {\texttt{ui\_anim}\\属性插值};
\node[modW] (uilyr)  at (8.2, -2.0) {\texttt{ui\_layer}\\8 按钮面板};
\node[modW] (status) at (8.2, -3.1) {\texttt{status\_line}\\8 行状态栏};
\node[modW] (txtbuf) at (8.2, -4.2) {\texttt{text\_buf}};
\node[modW] (font)   at (8.2, -5.3) {\texttt{font\_rom}};

% ---- HDMI ----
\node[modW] (hdmitop)at (8.2, -7.0) {\texttt{hdmi\_top}};
\node[modW] (enc)    at (8.2, -8.1) {\texttt{tmds\_enc} $\times$3};
\node[modW] (oser)   at (8.2, -9.2) {\texttt{tmds\_out}\\OSERDESE2 5:1};

% ---- SD 卡 ----
\node[modW] (sdtop)  at (8.2, -10.8){\texttt{sd\_ctrl\_top}};
\node[modW] (sdinit) at (8.2, -11.8){\texttt{sd\_init}\\CMD0/8/55/41};
\node[modW] (sdread) at (8.2, -12.8){\texttt{sd\_read}\\CMD18 多块};
\node[modW] (imgld)  at (4.2, -12.8){\texttt{sd\_img\_load}\\双写口};
\node[mem]   (fblcd) at (1.9, -12.6){帧缓冲\\480$\times$272\\RGB565};
\node[mem]   (fbhdmi)at (1.9, -14.4){帧缓冲\\320$\times$180\\RGB565};

\begin{scope}[on background layer]
\node[grpO, fit=(timing)(bg)(polar)(polrom)(mix)(rbrom)(demo)(wfilt),
      label={[grplbl]above:\texttt{disp\_top} $\times$2（参数门控 WITH\_UI/WITH\_TEXT）}] (gDisp) {};
\node[grpO, fit=(anim)(uilyr)(status)(txtbuf)(font),
      label={[grplbl]above:UI 层}] (gUi) {};
\node[grpG, fit=(hdmitop)(enc)(oser),
      label={[grplbl]above:HDMI (clk\_hdmi / clk\_5x)}] (gHdmi) {};
\node[grpP, fit=(sdtop)(sdinit)(sdread)(imgld)(fblcd)(fbhdmi),
      label={[grplbl]above:SD 卡与帧缓冲}] (gSd) {};
\node[draw=black!80, line width=1.0pt, rounded corners=6pt,
      fit=(gDisp)(gUi)(gHdmi)(gSd), inner sep=8pt,
      label={[font=\zihao{-5}\bfseries]above:FPGA —— XC7A100T-2FGG484I（显示与存储）}] (fpga) {};
\end{scope}

\draw[arBus] (gDisp.east) -- ++(0.4,0) |- node[lbl,pos=0.2,above]{RGB888+HS/VS/CLK} (lcd.east);
\draw[arBus] (gHdmi.east) -- ++(0.4,0) |- node[lbl,pos=0.2,above]{TMDS 4 对差分} (hdmi.east);
\draw[arD]   (sdcard.east) -- node[lbl,above]{SPI} ([yshift=0.4cm]gSd.west);
\draw[arBus] (fblcd.east) -- node[lbl,above]{LCD 全图} ([yshift=-1.4cm]gDisp.west);
\draw[arBus] (fbhdmi.east) -- ++(0.5,0) |- node[lbl,pos=0.3,above]{裁剪 2$\times$} ([yshift=-3.0cm]gDisp.west);
\draw[arBus] (gUi.south) -- ++(0,-0.35) -| node[lbl,pos=0.25,above]{ui\_sel / 配置} ([yshift=-2.6cm]gDisp.north -| gUi.south);
\draw[ar, dashed, draw=orange!70] ([yshift=0.9cm]fpga.north -| timing) -- node[lbl,above]{\texttt{clk\_pix} 12.5\,MHz} ([yshift=1.0cm]gDisp.north -| timing);
\draw[ar, dashed, draw=green!55!black] ([yshift=0.9cm]fpga.north -| hdmitop) -- node[lbl,above]{\texttt{clk\_5x} 125\,MHz} ([yshift=1.0cm]gHdmi.north -| hdmitop);

\end{tikzpicture}%
}
\caption{模块级框图（二）：双路显示（LCD 交互界面 / HDMI 展示大屏）与 SD 卡图片通道}
\label{fig:arch-b}
\end{figure}
\end{landscape}


\subsection{时钟架构}

\begin{table}[H]
\centering
\caption{时钟域划分}\label{tab:clk}
\begin{tabular}{llll}
\toprule[1.5pt]
\rowcolor{gray!25} 时钟 & 频率 & 来源 & 承载逻辑 \\
\midrule[1pt]
\texttt{clk\_sys}  & 48 MHz   & MMCM1（VCO 900\,\,÷\,18.75） & 音频、EQ、FFT、频谱、控制 \\
\texttt{clk\_pix}  & 12.5 MHz & MMCM1（900\,÷\,72）          & LCD 扫描 \\
\texttt{clk\_hdmi} & 25 MHz   & MMCM2（VCO 1000\,÷\,40）     & HDMI 像素 \\
\texttt{clk\_5x}   & 125 MHz  & MMCM2（1000\,÷\,8）          & TMDS 5:1 串行化 \\
BCLK               & 3.072 MHz& \textbf{外部}（WM8960）      & 不当时钟用，只过采样 \\
\bottomrule[1.5pt]
\end{tabular}
\end{table}

\textbf{关键决策}：BCLK/LRCLK 不接全局时钟网络，而是当作数据信号在
\texttt{clk\_sys} 域过采样出边沿选通。这样少一个时钟域、避开时钟脚抖动，
代价仅是 4 倍过采样（实际 15.6 倍，余量充足）。

%%%%%%%%%%%%%%%%%%%%%%%%%%%%%%%%%%%%%%%%%%%%%%%%%%%%%%%%%%%%%%%%%%%%%%%%%%%%%%%
\section{软硬件划分依据与接口设计}
%%%%%%%%%%%%%%%%%%%%%%%%%%%%%%%%%%%%%%%%%%%%%%%%%%%%%%%%%%%%%%%%%%%%%%%%%%%%%%%

\subsection{划分原则}

划分依据是“\textbf{什么是确定的、什么是易变的}”。

\begin{table}[H]
\centering
\caption{软硬件划分判据}\label{tab:split}
\begin{tabular}{p{6.2cm}p{3.6cm}p{3.6cm}}
\toprule[1.5pt]
\rowcolor{gray!25} 判据 & 硬件（RTL） & 软件（PC 脚本） \\
\midrule[1pt]
时序必须确定       & \checkmark & \\
每次运行结果都一样 & \checkmark & \\
需要人机交互调参   & \checkmark & \\
一次性的数据准备   &            & \checkmark \\
需要浮点数学       &            & \checkmark \\
需要第三方库       &            & \checkmark \\
\bottomrule[1.5pt]
\end{tabular}
\end{table}

\subsection{接口设计}

\begin{table}[H]
\centering
\caption{板级接口}\label{tab:if}
\begin{tabular}{p{2.4cm}p{1.2cm}p{4.6cm}p{5.0cm}}
\toprule[1.5pt]
\rowcolor{gray!25} 接口 & 方向 & 协议 / 电平 & 关键约束 \\
\midrule[1pt]
WM8960  & 双向 & I2C + I2S & WM8960 为 I2S 主，FPGA 为从 \\
LCD     & 输出 & RGB888 + HS/VS/CLK & 七项守恒检查 \\
HDMI    & 输出 & TMDS 4 对差分 & \textbf{板上无发送芯片}，FPGA 直连 \\
SD 卡   & 双向 & SPI & LVCMOS33，Bank 14 \\
UART    & 双向 & 115200 8N1 & ASCII 命令通道 \\
按键    & 输入 & 板载 + XPT2046 & 20 ms 消抖 + 300 ms 双击窗 \\
QSPI    & 双向 & W25Q128，4 线 & 掉电自启动 \\
\bottomrule[1.5pt]
\end{tabular}
\end{table}

\textbf{接口契约示例（帧缓冲写口）}：
\begin{lstlisting}[language=Verilog, basicstyle=\ttfamily\zihao{6},
                   frame=single, rulecolor=\color{boxFrame}]
output wire [18:0] wr_addr;   // 第 1 份：线性地址 = y*480 + x
output wire [15:0] wr_data;   //         RGB565
output wire        wr2_en;    // 第 2 份：只在裁剪区内有效
output wire [18:0] wr2_addr;  //         = (y-46)*320 + (x-80)
\end{lstlisting}

\textbf{设计约束}：\texttt{wr\_en} / \texttt{wr\_addr} / \texttt{wr2\_en} /
\texttt{wr2\_addr} 驱动 BRAM 的 \texttt{ENARDEN} / \texttt{ADDRARDADDR}，
\textbf{绝不能带异步复位}（否则触发 \texttt{DRC REQP-1839}）。

%%%%%%%%%%%%%%%%%%%%%%%%%%%%%%%%%%%%%%%%%%%%%%%%%%%%%%%%%%%%%%%%%%%%%%%%%%%%%%%
\section{优化过程与性能资源对比}
%%%%%%%%%%%%%%%%%%%%%%%%%%%%%%%%%%%%%%%%%%%%%%%%%%%%%%%%%%%%%%%%%%%%%%%%%%%%%%%

\subsection{优化总览}

\begin{table}[H]
\centering
\caption{五个主要优化阶段的量化收益}\label{tab:opt}
\begin{tabular}{cp{6.0cm}p{7.2cm}}
\toprule[1.5pt]
\rowcolor{gray!25} 阶段 & 优化内容 & 关键指标变化 \\
\midrule[1pt]
① & 同步复位替代异步复位 & \texttt{DPOR-1}: 11$\rightarrow$0；WNS: +1.117$\rightarrow$+2.156\,ns；LUT: 6802$\rightarrow$5648 \\
② & FFT 全并行 $\rightarrow$ 单蝶形时分复用 & DSP: 2048$\rightarrow$4（1/512）；BRAM: 0$\rightarrow$2 \\
③ & 新增 HDMI 大屏（80 柱 + 波形） & 新增功能，WNS 仍为正 \\
④ & SD 图片背景（双帧缓冲方案） & BRAM: 68.5$\rightarrow$100.5（+32 而非 +84） \\
⑤ & 综合警告清零 & 警告: 51$\rightarrow$0 \\
\bottomrule[1.5pt]
\end{tabular}
\end{table}

\subsection{优化 1：同步复位替代异步复位}

\textbf{现象}：\texttt{DPOR-1}$\times$11、\texttt{REQP-1839}$\times$180，
WNS 仅 +1.117\,ns（勉强收敛）。

\textbf{根因}：6 个文件用了 \texttt{always @(posedge clk or negedge rst\_n)}。
而本工程的 \texttt{rst\_n} 来自 \texttt{rst\_sync}（已是域内同步释放），
异步复位\textbf{纯代价}：DSP48 内部寄存器只支持同步复位，带异步复位就无法被吸收，
额外消耗约 1154 个 LUT 做置位/复位逻辑。

\textbf{量化收益}：

\begin{table}[H]
\centering
\caption{优化 1 前后对比}\label{tab:opt1}
\begin{tabular}{lccc}
\toprule[1.5pt]
\rowcolor{gray!25} 指标 & 优化前 & 优化后 & 变化 \\
\midrule[1pt]
\texttt{DPOR-1} & 11 & \textbf{0} & \checkmark \\
WNS            & +1.117 ns & \textbf{+2.156 ns} & +1.04 ns \\
Slice LUTs     & 6\,802 & \textbf{5\,648} & $-1\,154$ \\
\bottomrule[1.5pt]
\end{tabular}
\end{table}

\textbf{教训}：“已修复”必须按根因\textbf{全量重扫}——该根因第 4 次复发，
前 3 次都只改了报出来的那一处。

\subsection{优化 2：FFT 架构选型}

\begin{table}[H]
\centering
\caption{FFT 架构方案对比（1024 点，48 kHz 采样）}\label{tab:fft}
\begin{tabular}{lccccc}
\toprule[1.5pt]
\rowcolor{gray!25} 方案 & 蝶形单元 & DSP48 & BRAM & 整帧耗时 & 占帧周期 \\
\midrule[1pt]
全并行 & 512 & 2048 & 0 & 1 拍 & 0.04\% \\
4 路并行 & 128 & 512 & 0 & 7\,680 拍 & 0.3\% \\
\textbf{单蝶形（采用）} & \textbf{1} & \textbf{4} & \textbf{2} & \textbf{30\,720 拍} & \textbf{1.3\%} \\
\bottomrule[1.5pt]
\end{tabular}
\end{table}

\textbf{为什么够用}：帧周期 46.9\,ms，整帧耗时 $30\,720\times20.83\,\text{ns}=0.64$\,ms，
仅占 1.3\%，余量 75 倍。地址发生器用位操作而非除法：
$p=((c\gg s)\ll(s+1))\,|\,(c\wedge\text{mask})$，$q=p\,|\,(1\ll s)$。

实测：\textbf{1004 LUT / 126 FF / 2 BRAM36 / 4 DSP48}，WNS +2.54\,ns @ 48\,MHz。

\subsection{优化 3：SD 图片背景——四个方案的取舍}

\begin{table}[H]
\centering
\caption{SD 图片背景的四种候选方案}\label{tab:fb}
\begin{tabular}{cp{3.0cm}ccccc}
\toprule[1.5pt]
\rowcolor{gray!25} 方案 & 存储 & 每份 & 份数 & 总 BRAM & WNS & 结论 \\
\midrule[1pt]
A & 640$\times$360$\times$16 & 103 & 2 & 206 & — & 超器件 \\
B & 640$\times$360$\times$8  & 52  & 2 & 134 & +0.011 & 色带过重，实测否定 \\
C & 320$\times$180$\times$16 & 26  & 1 & 36.5 & +1.84 & 省资源但 LCD 有黑边 \\
\textbf{D} & \textbf{480$\times$272$\times$16 + 320$\times$180$\times$16} & \textbf{58+26} & \textbf{2} & \textbf{100.5} & \textbf{+1.78} & \textbf{采用} \\
\bottomrule[1.5pt]
\end{tabular}
\end{table}

\textbf{裁剪地址零 DSP 实现}：
\begin{lstlisting}[language=Verilog, basicstyle=\ttfamily\zihao{6},
                   frame=single, rulecolor=\color{boxFrame}]
crop_addr = (dy << 8) + (dy << 6) + dx;   // dy*256 + dy*64 = dy*320
\end{lstlisting}

\textbf{2 倍放大零成本}：
\begin{lstlisting}[language=Verilog, basicstyle=\ttfamily\zihao{6},
                   frame=single, rulecolor=\color{boxFrame}]
fb_addr = ((fb_ly >> 1) * FB_SW_U) + (fb_lx >> 1) + 1;   // 右移 1 位 = 除 2
\end{lstlisting}

\subsection{优化 4：综合警告清零（三类判据）}

\begin{table}[H]
\centering
\caption{51 条警告的分类处理}\label{tab:warn}
\begin{tabular}{p{2.2cm}c p{3.2cm} p{6.0cm}}
\toprule[1.5pt]
\rowcolor{gray!25} 类型 & 数量 & 处理 & 判据 \\
\midrule[1pt]
\textbf{真缺陷}   & \textbf{40} & \textbf{改设计} & \texttt{REQP-1839}：\texttt{wr2\_en}/\texttt{wr2\_addr} 驱动 BRAM 控制脚却带异步复位\\
前提不成立 & 50 & 豁免（写明判据） & \texttt{8-7129}：参数化模块高位地址超出实际使用范围 \\
假报警     & 1  & 豁免（附证据）   & \texttt{Route 35-328}：路由估计未命中，实测 WNS 全为正 \\
\bottomrule[1.5pt]
\end{tabular}
\end{table}

\subsection{最终资源与端到端延迟}

\begin{table}[H]
\centering
\caption{最终资源与性能}\label{tab:final}
\begin{tabular}{lcc}
\toprule[1.5pt]
\rowcolor{gray!25} 指标 & 值 & 说明 \\
\midrule[1pt]
Slice LUTs & 11\,716 & 18.48\% \\
Block RAM  & 100.5   & 74.44\% \\
DSP48      & 26      & 10.83\% \\
\texttt{clk\_sys} WNS  & \textbf{+1.779 ns} & 48 MHz，满足 \\
\texttt{clk\_hdmi} WNS & \textbf{+4.028 ns} & 25 MHz，满足 \\
\texttt{clk\_pix} WNS  & 远大于 0 & 12.5 MHz，满足 \\
构建警告   & \textbf{0 CRITICAL / 0 WARNING} & \\
仿真回归   & \textbf{41 / 41 PASS} & 全部逐位对拍 \\
\bottomrule[1.5pt]
\end{tabular}
\end{table}

\textbf{端到端延迟预算}：ADC$\rightarrow$I2S（20.83\,$\mu$s）$\rightarrow$
EQ（0.63\,$\mu$s）$\rightarrow$ FFT（0.64\,ms，占帧周期 1.3\%）$\rightarrow$
跨域（2\textasciitilde 3 拍）$\rightarrow$ LCD 扫描一帧（80\,ms）
$\Rightarrow$ \textbf{合计约 82\,ms，其中 80\,ms 是不可省的一帧扫描时间，
处理链路本身贡献不到 2\,ms}。

%%%%%%%%%%%%%%%%%%%%%%%%%%%%%%%%%%%%%%%%%%%%%%%%%%%%%%%%%%%%%%%%%%%%%%%%%%%%%%%
\section{大模型协作记录}
%%%%%%%%%%%%%%%%%%%%%%%%%%%%%%%%%%%%%%%%%%%%%%%%%%%%%%%%%%%%%%%%%%%%%%%%%%%%%%%

\subsection{所用工具}

\begin{table}[H]
\centering
\caption{AI 工具与用途}\label{tab:ai-tool}
\begin{tabular}{cp{3.6cm}p{3.2cm}p{5.6cm}}
\toprule[1.5pt]
\rowcolor{gray!25} 序号 & 工具 & 类别 & 主要用途 \\
\midrule[1pt]
1 & pi（编码代理 CLI）  & 智能体框架 & 工具调用、长会话、文件读写 \\
2 & Claude（Anthropic） & 主对话模型 & 设计讨论、RTL 编写、调试、文档 \\
3 & exa（MCP）          & 联网搜索   & 器件手册、协议细节、第三方实现比对 \\
4 & context7（MCP）     & 库文档检索 & Vivado Tcl / Verilator 参数 \\
5 & searchcode（MCP）   & 仓库搜索   & 找 TMDS 编码器独立实现做交叉验证 \\
\bottomrule[1.5pt]
\end{tabular}
\end{table}

大模型通过 \texttt{bash} 工具直接执行 \texttt{iverilog} / \texttt{verilator} /
\texttt{verible} / \texttt{vivado -mode batch}，形成
“改代码 $\rightarrow$ 跑仿真 $\rightarrow$ 看报告 $\rightarrow$ 再改”的闭环。

\subsection{使用环节}

\begin{table}[H]
\centering
\caption{使用环节与具体目的}\label{tab:ai-step}
\begin{tabular}{cp{3.0cm}p{8.8cm}}
\toprule[1.5pt]
\rowcolor{gray!25} 编号 & 环节 & 具体用途 \\
\midrule[1pt]
A-01 & 定点格式设计 & Q 格式、位宽预算、舍入与饱和策略 \\
A-02 & 滤波器结构选择 & 一阶低架 vs 二阶 peaking、级联顺序、条件数 \\
A-03 & FFT 架构权衡 & 存储式 vs 流式、单蝶形 vs 并行 \\
A-04 & 跨时钟域方案 & 格雷码 FIFO vs 双打 vs 电平翻转握手 \\
A-05 & 显示架构 & 扫描式渲染 vs framebuffer，参数化边界 \\
\midrule[0.5pt]
B-01 & RTL 编写 & 60 个 Verilog 模块 + 41 个 testbench \\
B-02 & 黄金模型 & Python 定点模型，逐位对拍 \\
B-03 & 生成器脚本 & 旋转因子、频谱分组表、字模、WM8960 寄存器表 \\
B-04 & 构建流程 & \texttt{build.tcl} / \texttt{program.tcl} / \texttt{run\_iv.sh} \\
B-05 & 调试 & 读波形、读综合报告、定位时序/资源/DRC 问题 \\
B-06 & 文档 & 16 个设计文档 + 825 条问题账本 \\
\bottomrule[1.5pt]
\end{tabular}
\end{table}

\subsection{主要提示方式}

\subsubsection{方式一：给“判据”而不是“要求”}

\textbf{低效提示}：“帮我写一个 SD 卡初始化模块。”

\textbf{高效提示}：
\begin{tcolorbox}[title=\textbf{用户}, colback=blue!2!white,
                  colframe=blue!30!black, fonttitle=\bfseries, breakable]
实现 SPI 模式的 SD 卡初始化：CMD0 $\rightarrow$ CMD8 $\rightarrow$
循环 (CMD55 + ACMD41) 直到 R1 的 bit0=0。\\
\textbf{判据}：① 第一次 \texttt{Go} 必须走完 5 步；
② \textbf{第二次 \texttt{Go} 也必须能重新走完}；
③ 初始化阶段卡不校验 CRC，用固定值即可。\\
\textbf{不要}：不要用状态机隐式等待，每步都要有超时计数。
\end{tcolorbox}

\textbf{结果}：一次生成即通过 5/5 判据，其中“第二次 \texttt{Go}”
这个判据后来救了大命（见下文例 3 的前身）。

\subsubsection{方式二：要求“交叉验证”而不是“相信”}

\textbf{提示}：TMDS 编码器我担心记错。请找 2\textasciitilde 3 份独立实现
（Xilinx 官方示例、开源核、论文附录），\textbf{逐行比对}，列出差异点再写；
若三者不一致就停下来告诉我。

\textbf{结果}：找到 3 份实现，其中 2 份在“DC 平衡”分支的 \texttt{cnt}
更新条件上写法不同但等价；确认后写入 \texttt{tmds\_enc.v}，一次通过。

\subsubsection{方式三：要求“量化”而不是“看起来对”}

\textbf{提示}：不要给“看起来合理”的结论。要求 ① 实际占用/总量/百分比；
② 与优化前对比；③ 变差了要指出是哪一项；
④ 时序报告里 WNS 是正是负、最差路径在哪。

\textbf{结果}：每次构建后都有一张可核对的指标表（见表~\ref{tab:final}）。

\subsection{自我纠错轨迹（五个典型案例）}

\subsubsection{例 1：\texttt{str.replace()} 静默成功}

\textbf{现象}：改完 \texttt{top.v} 的 \texttt{disp\_top} 实例参数，
脚本输出 OK，但上板实测 LCD 仍居中显示、颜色不对、画面错位重叠。

\textbf{定位}：
\begin{lstlisting}[language=bash, basicstyle=\ttfamily\zihao{6},
                   frame=single, rulecolor=\color{boxFrame}]
$ grep -n "FB_SW(" rtl/top.v
876:  .FB_SW(320), .FB_SH(180), .FB_W(320), .FB_H(180),   # 根本没改到
\end{lstlisting}

\textbf{根因}：\texttt{str.replace(old,new)} 在 \texttt{old} 不匹配时
\textbf{静默返回原串}，不报错不警告。

\textbf{修复}：\texttt{assert old in s, "搜索串没找到！"}——不匹配就崩，绝不静默。

\subsubsection{例 2：BRAM 被推断成 LUTRAM}

\textbf{现象}：合成报 \texttt{DRC UTLZ-1}：需 86\,452 个 \texttt{RAMD64E}，
只有 19\,000 个。

\textbf{定位链}：
① 直觉是“640$\times$360$\times$16 太大”，但算下来只要 103 块（器件有 135 块）；
② 看报告：\textbf{“LUT as Distributed RAM”} 超限，说明根本没映射到 BRAM；
③ 检查：\texttt{mem} 被两个不同时钟读；
④ 加 \texttt{(* ram\_style = "block" *)}，重合成\textbf{仍然 LUTRAM}；
⑤ 查 UG473：BRAM 推断模板只支持 1 写 1 读或真双口（需原语）。

\textbf{修复}：改成两份不同尺寸的缓冲实例。

\subsubsection{例 3：\texttt{wr\_en} 类信号误放进异步复位块}

\textbf{现象}：综合报 40 条 \texttt{REQP-1839}。

\textbf{定位}：检查 \texttt{sd\_img\_load.v}——我给第二个写口加
\texttt{wr2\_en}/\texttt{wr2\_addr} 时，把它们放进了 \texttt{if (!rst\_n)} 块。
而\textbf{同一个文件里 \texttt{wr\_addr}/\texttt{wr\_en} 的注释明确写着“不能带异步复位”}。

\textbf{修复}：移出异步复位块，改由 \texttt{start} 清零 + 每拍给默认值。

\subsubsection{例 4：\texttt{y[8:3]} 在 360 行时回绕}

\textbf{现象}：LCD 与 HDMI 的程序化背景不一致。

\textbf{定位}：\texttt{bg\_src} 是纯组合函数，同 $(x,y)$ 必同色；但两实例
\texttt{VDISP} 不同（272 / 360）。渐变代码 \texttt{wire [5:0] t = y[8:3];}
——272 行时 $y[8:3]=0\ldots33$ 正确；360 行时 $y\ge256$ 使
\textbf{$y[8:3]$ 回绕到 0}，下半屏渐变从头开始。

\textbf{修复}：\texttt{YSH = (VDISP>=512) ? 4 : 3}，\texttt{t = y[YSH +: 6]}。

\subsubsection{例 5：颜色不对——数清“反转层次”}

\textbf{现象}：HDMI 输出图片，原图深蓝却显示成暖黄；LCD 正常。

\textbf{定位链}：
① “深蓝$\rightarrow$暖黄”是 \textbf{R/B 互换的教科书特征}；
② 数据链只有一处可能：\texttt{disp\_top} 的 RGB565$\rightarrow$RGB888 展开，
高字节接 \texttt{fb\_rgb[4:0]}（B）、低字节接 \texttt{fb\_rgb[15:11]}（R）；
③ 但这行是按 LCD 实测写的，且 LCD 显示正确；
④ 推理：\textbf{LCD 面板接线本身是 BGR}，这里反一次、面板又反一次，抵消；
⑤ \textbf{HDMI 走标准 RGB，没有第二次反}，所以必须在这里换回来。

\textbf{修复}：仅 HDMI 路径加
\texttt{.pix\_rgb(\{r[7:0], r[15:8], r[23:16]\})}。

\subsection{采纳、人工修改与核验}

\begin{table}[H]
\centering
\caption{AI 输出的采纳情况}\label{tab:adopt}
\begin{tabular}{cp{2.2cm}p{5.0cm}p{3.4cm}}
\toprule[1.5pt]
\rowcolor{gray!25} 编号 & 采纳 & 人工修改 & 核验结果 \\
\midrule[1pt]
A-01 & 采纳     & 无 & \checkmark 逐位一致 \\
A-02 & 部分采纳 & 结构从 4 段改 5 段，加 8 kHz peaking & \checkmark 频响误差 $<$0.5 dB \\
A-03 & 采纳     & 无 & \checkmark 2048 频点逐位一致 \\
A-04 & 部分采纳 & 格雷码 FIFO 改为“电平翻转+双打+边沿检测” & \checkmark 4/4 PASS \\
A-05 & 部分采纳 & 加 \texttt{WITH\_UI}/\texttt{WITH\_TEXT}/\texttt{FB\_*} 参数族 & \checkmark 130\,560 像素逐点一致 \\
B-01 & 采纳     & 多次人工纠错（见 5.4） & \checkmark 41/41 \\
B-02 & 采纳     & 无 & \checkmark 逐位一致 \\
B-03 & 采纳     & \texttt{gen\_spectrum\_map.py} 后参数化 & \checkmark 60/80 均验证 \\
B-04 & 部分采纳 & Flash 烧写改成“先 \texttt{program\_hw\_devices} 建 SPI 通道” & \checkmark 43 秒成功 \\
B-05 & 采纳     & 无 & \checkmark WNS 全为正 \\
B-06 & 采纳     & 全部数字人工核对 & \checkmark 与报告一致 \\
\bottomrule[1.5pt]
\end{tabular}
\end{table}

\textbf{核验方法论}：① 逐位对拍（不允许“看起来差不多”）；
② 独立真值（\texttt{numpy.fft} / \texttt{\$atan2}）；
③ 穷举（512 个 bin 分组、5132 组地址、49\,920 个像素）；
④ 总量守恒（一帧 150\,150 拍、130\,560 个有效像素）；
⑤ 读回验证（帧缓冲逐点写入再读回）。

%%%%%%%%%%%%%%%%%%%%%%%%%%%%%%%%%%%%%%%%%%%%%%%%%%%%%%%%%%%%%%%%%%%%%%%%%%%%%%%
\section{技能包的提炼过程}\label{sec:skill}
%%%%%%%%%%%%%%%%%%%%%%%%%%%%%%%%%%%%%%%%%%%%%%%%%%%%%%%%%%%%%%%%%%%%%%%%%%%%%%%

\subsection{三级沉淀体系}

\begin{table}[H]
\centering
\caption{三级知识沉淀}\label{tab:levels}
\begin{tabular}{p{2.6cm}p{4.4cm}cp{4.6cm}}
\toprule[1.5pt]
\rowcolor{gray!25} 层级 & 载体 & 规模 & 作用 \\
\midrule[1pt]
问题账本 & \texttt{docs/08-skills-index.md} & \textbf{825 条} & 现象/根因/修复/量化改善，原始可检索 \\
技能包   & \texttt{.pi/skills/*/SKILL.md}   & \textbf{4 个} & 触发条件/方法/判据/边界，可复用 \\
状态快照 & \texttt{docs/14-status-snapshot.md} & 1 份 & 一眼看清“能用的/受限的/待办的” \\
\bottomrule[1.5pt]
\end{tabular}
\end{table}

\subsection{四个技能包}

\begin{table}[H]
\centering
\caption{技能包来源、验证与边界}\label{tab:skills}
\begin{tabular}{p{2.4cm}p{3.6cm}p{3.4cm}p{3.6cm}}
\toprule[1.5pt]
\rowcolor{gray!25} 技能包 & 来源失败 & 如何验证有效 & 边界 \\
\midrule[1pt]
\texttt{fixed-point-dsp} & EQ 频响差 3 dB（位宽不足）；\texttt{WAVE\_AMP} 改大反而变小（静默截断）；某频点发散（条件数） & 重做 EQ，频响误差 3 dB$\rightarrow$\textbf{0.5 dB}；截断问题后续 6 个参数\textbf{零复发} & 适用：定点 IIR/FIR/FFT，16\textasciitilde 32 bit。不适用：浮点、位宽 $<$12 bit \\
\texttt{fpga-rtl-flow} & 找不到 Vivado；DRC 在 GUI 不报；时序违例方向错；\texttt{place} 失败伪装成综合错；构建 9m38s 太慢 & 构建 \textbf{9m38s$\rightarrow$2m57s}；回归 \textbf{20m54s$\rightarrow$1m30s}；后续 20+ 次构建\textbf{零次返工} & 适用：纯 RTL、无 block design。不适用：IP integrator、部分重配置 \\
\texttt{rtl-verification} & 在信号自己更新的沿上读它；非阻塞赋值读不到新值；跨帧/帧内状态共用复位；只测一次漏掉重配路径 & 41 个 TB 中\textbf{至少 12 个第一次跑就红}；账本里\textbf{不到 10 条是纯仿真能发现}，但那 10 条都是上板必花几天的 & 适用：确定性数字逻辑。不适用：模拟/混合信号 \\
\texttt{contributor} & 重复劳动、项目不活跃、未读 CONTRIBUTING & —— & 适用：活跃开源项目 \\
\bottomrule[1.5pt]
\end{tabular}
\end{table}

\subsection{从“观察”到“技能”的抽象链路}

以 \texttt{rtl-verification} 为例：

\begin{tcolorbox}[colback=gray!4, colframe=boxFrame, breakable,
                  title=\textbf{抽象过程}]
\textbf{原始观察（账本 412）}：TB 里 \texttt{@(negedge sd\_clk)} 采样，
结果总差一拍。\\
$\downarrow$ 追问“为什么”\\
\textbf{根因}：\texttt{sd\_clk} 由 \texttt{div\_clk} 分频而来，
\texttt{negedge sd\_clk} 恰好等于 \texttt{posedge div\_clk}，
在“上升沿那一刻”读到的是旧值。\\
$\downarrow$ 再问“是个例还是通类”\\
\textbf{聚类（账本 413/414/415）}：非阻塞赋值在同一 \texttt{always} 读不到新值；
\texttt{acmd41\_cnt} 在 CS 沿被清零导致永远回不到第 3 次；
边际推理调不动时应换显式计数。\\
$\downarrow$ 抽象成规则\\
\textbf{技能包条目}：“永远不要在信号自己更新的沿上读它”；
“复位的作用域要精确界定：跨帧状态 $\neq$ 帧内状态”。\\
$\downarrow$ 写成可执行检查\\
\textbf{判据}：写 TB 时自问——这个 \texttt{@(posedge/negedge X)}
和 X 的驱动沿是同一个吗？是则改成 \texttt{@(posedge clk); \#1;} 或换一个沿。
\end{tcolorbox}

\subsection{边界如何确定}

\textbf{判据：每条技能都必须能回答三个问题}——
① 它解决什么现象？② 它的前提是什么？
③ 前提不成立时会怎样？

\textbf{反面例子（避免把特例当通例）}：
\begin{tcolorbox}[colback=red!3, colframe=red!40, breakable]
$\times$ \textbf{不要写}：“BRAM 一定要加 \texttt{(* ram\_style = "block" *)}”
——实测该属性在“1 写 2 读跨时钟”下\textbf{无效}。\\
$\checkmark$ \textbf{应该写}：“BRAM 推断由\textbf{访问模式}决定，属性只是建议。
1W1R 可靠；1W2R 跨时钟必退化成 LUTRAM，属性也救不了。”
\end{tcolorbox}

%%%%%%%%%%%%%%%%%%%%%%%%%%%%%%%%%%%%%%%%%%%%%%%%%%%%%%%%%%%%%%%%%%%%%%%%%%%%%%%
\section{复现说明}
%%%%%%%%%%%%%%%%%%%%%%%%%%%%%%%%%%%%%%%%%%%%%%%%%%%%%%%%%%%%%%%%%%%%%%%%%%%%%%%

\subsection{环境要求}

\begin{table}[H]
\centering
\caption{软件环境}\label{tab:env}
\begin{tabular}{lll}
\toprule[1.5pt]
\rowcolor{gray!25} 工具 & 版本 & 说明 \\
\midrule[1pt]
Vivado     & 2020.2 & 免费 WebPACK 版即可（7A100T 在支持列表内）\\
iverilog   & 12.0   & \texttt{pacman -S iverilog} / \texttt{apt install iverilog} \\
Verilator  & 5.052  & \texttt{pacman -S verilator} \\
Verible    & latest & \texttt{mason} 或 GitHub release \\
Python     & $\ge$3.10 & 系统自带 \\
Pillow     & $\ge$10.0 & 仅 \texttt{mkimage.py} 处理 PNG/JPG 时需要 \\
中文字体   & Noto CJK & 编译本文档需要 \\
\bottomrule[1.5pt]
\end{tabular}
\end{table}

\subsection{从零开始的完整步骤}

\subsubsection{步骤 1：获取代码}
\begin{lstlisting}[language=bash, basicstyle=\ttfamily\zihao{6},
                   frame=single, rulecolor=\color{boxFrame}]
git clone https://github.com/flyingyuya/Sonic-Speed-Parallelism.git
cd Sonic-Speed-Parallelism
\end{lstlisting}

\subsubsection{步骤 2：工具链自检}
\begin{lstlisting}[language=bash, basicstyle=\ttfamily\zihao{6},
                   frame=single, rulecolor=\color{boxFrame}]
bash scripts/sim/run_iv.sh tools
# 预期：[ok] iverilog / verilator / verible 三行
# 若 Vivado 找不到：export PATH=$HOME/Xilinx/Vivado/2020.2/bin:$PATH
\end{lstlisting}

\subsubsection{步骤 3：全量仿真（约 1 分 30 秒）}
\begin{lstlisting}[language=bash, basicstyle=\ttfamily\zihao{6},
                   frame=single, rulecolor=\color{boxFrame}]
bash scripts/sim/run_iv.sh
# 预期：通过 41  失败 0  跳过 0     <-- 这是环境是否正确的黄金标准
# 单项：bash scripts/sim/run_iv.sh fft_core
# 含 tb_top 的全量：FULL=1 bash scripts/sim/run_iv.sh
\end{lstlisting}

\subsubsection{步骤 4：静态检查}
\begin{lstlisting}[language=bash, basicstyle=\ttfamily\zihao{6},
                   frame=single, rulecolor=\color{boxFrame}]
verible-verilog-lint --rules_config_search $(find rtl sim -name '*.v')
# 预期：无输出（退出码 0）
\end{lstlisting}

\subsubsection{步骤 5：综合与实现（约 9\textasciitilde 10 分钟）}
\begin{lstlisting}[language=bash, basicstyle=\ttfamily\zihao{6},
                   frame=single, rulecolor=\color{boxFrame}]
vivado -mode batch -source scripts/vivado/build.tcl \
       -tclargs xc7a100tfgg484-2 top constrs
# 预期：构建 0  CRITICAL 0  WARNING 0
#       LUT 11716 (18.48%)  BRAM 100.5 (74.44%)
#       clk_sys +1.779ns    clk_hdmi +4.028ns
# 只综合（约 3 分钟，看资源够不够）：
SYNTH_ONLY=1 vivado -mode batch -source scripts/vivado/build.tcl \
       -tclargs xc7a100tfgg484-2 top constrs
\end{lstlisting}

\subsubsection{步骤 6：准备 SD 卡（需要 \texttt{sudo}）}
\begin{lstlisting}[language=bash, basicstyle=\ttfamily\zihao{6},
                   frame=single, rulecolor=\color{boxFrame}]
lsblk -o NAME,SIZE,TYPE,TRAN | grep -E "sd|mmc"     # 找到读卡器，如 sda

# 先 dry-run（强烈建议）
python3 scripts/sdcard/mkimage.py --dry-run /dev/sda img1.png img2.png img3.png
# 预期每张图 261120 字节 = 480x272x2  <-- 数字对了就说明格式正确

# 正式写入（脚本内建系统盘保护）
echo YES | sudo python3 scripts/sdcard/mkimage.py /dev/sda img1.png img2.png img3.png
\end{lstlisting}

\subsubsection{步骤 7：上板}
\begin{lstlisting}[language=bash, basicstyle=\ttfamily\zihao{6},
                   frame=single, rulecolor=\color{boxFrame}]
# 1) 验证 JTAG
vivado -mode batch -source scripts/vivado/program.tcl -tclargs detect
# 预期：发现 JTAG 目标 ... 认到目标器件 xc7a100t

# 2) 下载到 SRAM（10 秒，掉电即失，开发期用这个）
vivado -mode batch -source scripts/vivado/program.tcl

# 3) 固化到 Flash（43 秒，掉电保持，定型后才用）
vivado -mode batch -source scripts/vivado/program.tcl -tclargs flash
# 预期：Mfg ID ef / Erase-Program-Verify successful
# 注意：JTAG 优先于 QSPI，验证 Flash 启动必须拔掉 JTAG 线再上电
\end{lstlisting}

\subsection{验收清单}

\begin{table}[H]
\centering
\caption{上板验收清单}\label{tab:accept}
\begin{tabular}{cp{6.4cm}p{6.0cm}}
\toprule[1.5pt]
\rowcolor{gray!25} \# & 检查项 & 预期 \\
\midrule[1pt]
1  & \texttt{run\_iv.sh} & \textbf{41/41 PASS} \\
2  & \texttt{verible}    & 0 违规 \\
3  & 构建               & \textbf{0 CRITICAL / 0 WARNING} \\
4  & \texttt{led[1]} 闪烁 & FPGA 活着，时钟/复位正常 \\
5  & \texttt{led[0]} 常亮 & WM8960 I2C 配置完成 \\
6  & LCD 显示           & 网格 + 渐变背景 \\
7  & 按 \texttt{BG} 键  & 在程序化背景 / SD 图片间切换 \\
8  & 按 \texttt{SD} 键  & 3 张图片循环 \\
9  & 对麦克风说话       & 频谱柱跳动、波形变化 \\
10 & 按 \texttt{SPK} 键 & 喇叭静音/取消 \\
11 & HDMI 显示器        & 640$\times$360 有画面（80 柱 + 波形 + 图片）\\
12 & UART 115200        & 上电 50 ms 打横幅；发 \texttt{V} 切视图 \\
\bottomrule[1.5pt]
\end{tabular}
\end{table}

\subsection{常见问题}

\begin{table}[H]
\centering
\caption{常见问题与解决}\label{tab:faq}
\begin{tabular}{p{4.0cm}p{3.8cm}p{5.4cm}}
\toprule[1.5pt]
\rowcolor{gray!25} 现象 & 原因 & 解决 \\
\midrule[1pt]
\texttt{run\_iv.sh} 全失败 & iverilog 太老 & 需 $\ge$11.0 \\
综合报 \texttt{UTLZ-1}（LUTRAM 超限） & 存储器多端口访问模式不匹配推断模板 & 复制实例，不要加读口 \\
\texttt{place\_design} 失败但无时序错 & BRAM $>$95\% 导致布线拥塞 & 看 \texttt{post\_route\_util.rpt}，降位深/降分辨率 \\
上板 LCD 全白/全黑 & LCD 时序参数错 & 跑 \texttt{run\_iv.sh lcd\_timing} \\
SD 图片乱码/斜切 & 卡里数据尺寸与 RTL 的 W/H 不匹配 & 重新写卡，确认每图 261\,120 字节 \\
HDMI 颜色偏黄 & R/B 通道互换 & 见 5.4 例 5 \\
HDMI 复位后偶发右半屏撕裂 & \textbf{已知未解决}（TMDS 层） & 重按复位可恢复；需 ILA/示波器 \\
线路输入无声 & 板上 \texttt{C4}(1\,$\mu$F) 开路 & 需补焊（代码已就绪）\\
\bottomrule[1.5pt]
\end{tabular}
\end{table}

%%%%%%%%%%%%%%%%%%%%%%%%%%%%%%%%%%%%%%%%%%%%%%%%%%%%%%%%%%%%%%%%%%%%%%%%%%%%%%%
\section*{附录 A\quad 关键指标汇总}
\addcontentsline{toc}{section}{附录 A\quad 关键指标汇总}
%%%%%%%%%%%%%%%%%%%%%%%%%%%%%%%%%%%%%%%%%%%%%%%%%%%%%%%%%%%%%%%%%%%%%%%%%%%%%%%

\begin{table}[H]
\centering
\caption{关键指标}\label{tab:appendix}
\begin{tabular}{p{2.4cm}p{5.6cm}p{5.0cm}}
\toprule[1.5pt]
\rowcolor{gray!25} 类别 & 指标 & 值 \\
\midrule[1pt]
\multirow{3}{*}{规模} & RTL 文件 / 行数 & 60 / $\approx$12\,000 \\
                      & testbench 数量  & 41 \\
                      & 问题账本条目    & 825 \\
\midrule[0.5pt]
\multirow{3}{*}{资源} & Slice LUTs & 11\,716（18.48\%）\\
                      & Block RAM  & 100.5（74.44\%）\\
                      & DSP48      & 26（10.83\%）\\
\midrule[0.5pt]
\multirow{3}{*}{时序} & \texttt{clk\_sys} 48 MHz  & WNS +1.779 ns \\
                      & \texttt{clk\_hdmi} 25 MHz & WNS +4.028 ns \\
                      & \texttt{clk\_pix} 12.5 MHz& WNS 远大于 0 \\
\midrule[0.5pt]
\multirow{3}{*}{质量} & 构建警告 & 0 CRITICAL / 0 WARNING \\
                      & 静态检查 & Verible 0 / Verilator 0 \\
                      & 仿真回归 & 41/41 PASS \\
\midrule[0.5pt]
\multirow{3}{*}{性能} & 音频采样 & 48 kHz / 24 bit \\
                      & FFT 整帧 & 0.64 ms（占帧周期 1.3\%）\\
                      & 端到端延迟 & $\approx$82 ms（其中 80 ms 是 LCD 一帧扫描）\\
\midrule[0.5pt]
\multirow{4}{*}{功能} & LCD  & 480$\times$272，8 按钮 + 8 行状态栏 + 双背景 + SD 图片 \\
                      & HDMI & 640$\times$360，80 根独立柱 + 波形 + SD 图片（裁剪 2$\times$ 放大）\\
                      & 音频 & 话筒 / EQ 5 段 / 双声道 / 静音 / 回环 \\
                      & 存储 & SD（CMD18 多块读）/ QSPI Flash 掉电自启动 \\
\bottomrule[1.5pt]
\end{tabular}
\end{table}

%%%%%%%%%%%%%%%%%%%%%%%%%%%%%%%%%%%%%%%%%%%%%%%%%%%%%%%%%%%%%%%%%%%%%%%%%%%%%%%
\section*{附录 B\quad 与商业方案的对比}
\addcontentsline{toc}{section}{附录 B\quad 与商业方案的对比}
%%%%%%%%%%%%%%%%%%%%%%%%%%%%%%%%%%%%%%%%%%%%%%%%%%%%%%%%%%%%%%%%%%%%%%%%%%%%%%%

\begin{table}[H]
\centering
\caption{与商业方案对比}\label{tab:cmp}
\begin{tabular}{p{2.4cm}p{2.8cm}p{2.4cm}p{2.4cm}p{3.4cm}}
\toprule[1.5pt]
\rowcolor{gray!25} 维度 & PC 软件 & Zynq+PYNQ & 商用 ASIC & \textbf{本作品} \\
\midrule[1pt]
端到端延迟 & 10\textasciitilde 100 ms & 1\textasciitilde 10 ms & $<$1 ms & \textbf{$\approx$82 ms} \\
处理链路延迟 & 不确定 & 依赖 OS/缓存 & 确定 & \textbf{$<$2 ms，可精确到拍} \\
在线调参 & 灵活 & 灵活 & 不可 & \textbf{1 拍生效} \\
启动时间 & 秒级 & 秒级 & 即时 & \textbf{微秒级} \\
IP 依赖 & 无 & HLS+AXI & —— & \textbf{零} \\
可审计性 & 高 & 低（HLS 不可见） & 无 & \textbf{每一行可查} \\
成本 & 低 & 中 & 高 & \textbf{低} \\
\bottomrule[1.5pt]
\end{tabular}
\end{table}

\end{document}
